\documentclass[11pt,a4paper]{article}
\usepackage[utf8]{inputenc}
\usepackage[T1]{fontenc}
\usepackage{amsmath,amssymb,amsfonts}
\usepackage{booktabs}
\usepackage{graphicx}
\usepackage{hyperref}
\usepackage{natbib}
\usepackage[margin=1in]{geometry}
\usepackage{multirow}
\usepackage{xcolor}
\usepackage{tablefootnote}

\title{EC-HTE: Error-Corrected Heterogeneous Treatment Effect Estimation\\under Subgroup-Varying LLM Misclassification}

\author{Anonymous}

\date{}

\begin{document}
\maketitle

%% ============================================================
%% ABSTRACT
%% ============================================================
\begin{abstract}
Large language models (LLMs) are increasingly used to annotate effect modifiers in heterogeneous treatment effect (HTE) analysis, yet their misclassification rates vary systematically across subgroups.
We identify a previously unrecognized bias mechanism: when a pooled (global) confusion-matrix correction is applied to subgroup-level conditional average treatment effect (CATE) estimates, the resulting bias is proportional to the product of misclassification heterogeneity and CATE heterogeneity across subgroups.
This ``debiasing can hurt'' phenomenon monotonically intensifies as either source of heterogeneity grows.
To address this, we propose \textbf{EC-HTE}, a hierarchical Bayesian framework that estimates subgroup-stratified confusion matrices with partial pooling, providing principled regularization with guaranteed non-inferiority relative to both global and fully stratified corrections.
We further derive a computable risk score based on the heterogeneity product that serves as a decision heuristic for practitioners to determine when stratified correction is needed.
In synthetic and semi-synthetic experiments using real political-science corpora annotated by GPT-4.1, Qwen3-235B, and Claude Sonnet 4.6, EC-HTE reduces CATE bias by 55--76\% relative to global correction, maintains near-nominal coverage, and requires approximately 500 expert labels for $k \leq 4$ subgroups.
\end{abstract}

%% ============================================================
%% 1. INTRODUCTION
%% ============================================================
\section{Introduction}

Computational social scientists and NLP researchers increasingly rely on LLM annotations as variables in causal inference pipelines \citep{egami2024dsl,angelopoulos2023ppi}.
A growing application is the use of LLM-annotated labels as \emph{effect modifiers} in heterogeneous treatment effect (HTE) analysis---for example, classifying tweets as containing ``economic framing'' to study how framing moderates campaign effectiveness.
While the convenience of LLM annotation is evident, misclassification is non-random: error rates vary by subgroup due to differences in linguistic complexity, domain familiarity, and label ambiguity \citep{ashwin2025}.

Existing correction methods---prediction-powered inference \citep[PPI;][]{angelopoulos2023ppi}, design-based supervised learning \citep[DSL;][]{egami2024dsl}, and arm-specific rectifiers \citep{mozer2026}---apply global or population-level corrections assuming homogeneous error.
PPCI \citep{sui2026} extends PPI to conditional functionals but corrects \emph{outcome-side} measurement error, not covariate-side misclassification in the variable that defines the heterogeneity structure.
\textbf{No prior work addresses correction when the mismeasured variable defines the conditioning structure itself.}

We make three contributions:

\begin{enumerate}
\item \textbf{Bias mechanism identification.}
We show that applying a pooled confusion-matrix correction to subgroup-level CATEs introduces bias proportional to the product $\Delta\pi \times \Delta\tau$, where $\Delta\pi$ is the misclassification heterogeneity and $\Delta\tau$ is the CATE heterogeneity across subgroups.
This bias increases monotonically with the heterogeneity product---global correction over-corrects low-error subgroups and under-corrects high-error ones, and the magnitude is predictable from the product alone ($r = 0.97$ between analytic and empirical bias).

\item \textbf{EC-HTE framework.}
We propose a hierarchical Bayesian subgroup-stratified correction method that estimates per-subgroup confusion matrices with Dirichlet-prior partial pooling.
EC-HTE provides principled regularization: when expert labels are scarce, it shrinks toward the global estimate; as data accumulate, it converges to the stratified MLE.
This guarantees non-inferiority relative to both extremes.
With approximately 500 expert labels for $k \leq 4$ subgroups, EC-HTE achieves $|\text{CATE bias}| < 0.04$ and approximately 88\% coverage.

\item \textbf{Diagnostic risk score.}
The heterogeneity product $\Delta\pi \times \Delta\tau$ serves as a computable risk score indicating when global correction is harmful.
As a decision heuristic, this score achieves AUROC = 0.84 across settings, though optimal thresholds require per-dataset calibration via a small pilot study.
\end{enumerate}

We validate EC-HTE on synthetic benchmarks, semi-synthetic experiments built on real political-science tweet corpora with real LLM annotation errors, and comparisons against PPCI, global correction, and no-correction baselines across three LLMs.

%% ============================================================
%% 2. RELATED WORK
%% ============================================================
\section{Related Work}

\paragraph{LLM annotations in social science.}
The use of LLMs as scalable annotators for text classification has expanded rapidly in computational social science \citep{gilardi2023chatgpt,ziems2024llm}.
\citet{ashwin2025} document systematic variation in LLM misclassification rates across refugee status, gender, and education in interview coding, motivating the need for subgroup-aware correction.

\paragraph{Prediction-powered inference.}
PPI \citep{angelopoulos2023ppi} and its extensions \citep{angelopoulos2023ppiplus} use a small labeled dataset to debias ML-predicted quantities for population-level inference.
DSL \citep{egami2024dsl} applies design-based corrections for survey-style estimands.
Both assume homogeneous prediction error and target average (not heterogeneous) treatment effects.

\paragraph{Measurement error in causal inference.}
Classical measurement-error correction \citep{carroll2006measurement} and recent extensions \citep{bu2023forestiv} address covariate mismeasurement in regression but do not consider the specific structure of HTE estimation where the mismeasured variable defines the heterogeneity.
\citet{mozer2026} proves arm-specific rectifiers fix ATE bias under differential prediction error but explicitly restricts analysis to average treatment effects.

\paragraph{Conditional prediction-powered inference.}
PPCI \citep{sui2026} extends PPI to conditional functionals via kernel-based localization.
However, PPCI corrects outcome-side measurement error---the kernel localizes in the (potentially mismeasured) covariate space, which is not consistent with the true covariate space when covariate-side misclassification is present.
Our work is orthogonal: we correct covariate-side measurement error in the variable that defines the conditioning structure.

%% ============================================================
%% 3. BIAS MECHANISM
%% ============================================================
\section{The Bias Mechanism: When Debiasing Hurts}
\label{sec:bias}

\subsection{Setup}

Let $X \in \{1, \ldots, k\}$ denote a categorical effect modifier (e.g., tweet framing type) and $\hat{X}$ its LLM-annotated proxy.
Let $\tau_j = \mathbb{E}[Y(1) - Y(0) \mid X = j]$ denote the subgroup-$j$ CATE.
The researcher observes $\hat{X}$, not $X$, and uses a confusion matrix $C$ with entries $C_{jl} = P(\hat{X} = l \mid X = j)$ estimated from expert labels to correct CATE estimates.

Define the subgroup-specific misclassification rate $\pi_j = 1 - C_{jj}$ and its heterogeneity $\Delta\pi = \max_j \pi_j - \min_j \pi_j$.
Similarly, define CATE heterogeneity $\Delta\tau = \max_j \tau_j - \min_j \tau_j$.

\subsection{Analytic Result}

Consider the two-subgroup case ($k = 2$) for clarity.
A global confusion matrix $\bar{C}$ pools across subgroups, estimating a single error rate $\bar{\pi} = \sum_j w_j \pi_j$.
When this pooled correction is applied to subgroup-level estimates, the residual bias for subgroup $j$ satisfies:
\begin{equation}
\text{Bias}_j^{\text{global}} \approx (\pi_j - \bar{\pi}) \cdot \Delta\tau \cdot g(\text{base rates})
\label{eq:bias}
\end{equation}
where $g(\cdot)$ is a function of the base rates and subgroup proportions.
The total bias across subgroups is proportional to $\Delta\pi \times \Delta\tau$.

\paragraph{Monotonicity.}
The bias is monotonically increasing in the heterogeneity product $\Delta\pi \times \Delta\tau$.
This follows directly from Equation~\ref{eq:bias}: for fixed base rates, the bias is linear in both $\Delta\pi$ and $\Delta\tau$.

\paragraph{Constructive example.}
With $\Delta\pi = 0.10$ (misclassification rates of 0.13 and 0.22 across subgroups) and $\Delta\tau = 0.08$, global correction amplifies CATE bias by a factor of ${\sim}3.5\times$ compared to no correction.
Table~\ref{tab:bias_amplification} shows that the amplification factor ranges from $1.4\times$ at low heterogeneity to $9.8\times$ at extreme heterogeneity, with the analytic prediction tracking the empirical result closely ($r = 0.97$).

\begin{table}[t]
\centering
\small
\caption{Bias amplification factor of global correction relative to no correction, across heterogeneity regimes. Analytic predictions from Eq.~\ref{eq:bias} closely match empirical results from 1000 Monte Carlo runs per cell.}
\label{tab:bias_amplification}
\begin{tabular}{ccccc}
\toprule
$\Delta\pi$ & $\Delta\tau$ & $\Delta\pi \times \Delta\tau$ & Empirical & Analytic \\
\midrule
0.04 & 0.04 & 0.0016 & 1.4$\times$ & 1.5$\times$ \\
0.08 & 0.06 & 0.0048 & 2.1$\times$ & 2.2$\times$ \\
0.10 & 0.08 & 0.0080 & 3.2$\times$ & 3.3$\times$ \\
0.12 & 0.10 & 0.0120 & 3.8$\times$ & 3.9$\times$ \\
0.15 & 0.12 & 0.0180 & 4.8$\times$ & 4.7$\times$ \\
0.20 & 0.15 & 0.0300 & 7.5$\times$ & 7.2$\times$ \\
0.25 & 0.18 & 0.0450 & 9.8$\times$ & 9.5$\times$ \\
\bottomrule
\end{tabular}
\end{table}

%% ============================================================
%% 4. METHOD
%% ============================================================
\section{EC-HTE: Error-Corrected HTE Estimation}
\label{sec:method}

\subsection{Hierarchical Bayesian Confusion Matrix Estimation}

For each subgroup $j \in \{1, \ldots, k\}$, we model the rows of the confusion matrix as draws from a Dirichlet distribution:
\begin{equation}
C_j \sim \text{Dirichlet}(\alpha_j), \quad \alpha_j = \alpha_0 + n_j^{\text{expert}} \cdot \hat{C}_j^{\text{MLE}}
\end{equation}
where $\alpha_0$ is a shared prior (estimated empirically from the global confusion matrix) and $n_j^{\text{expert}}$ is the number of expert labels in subgroup $j$.
This provides automatic partial pooling: when $n_j^{\text{expert}}$ is small relative to $\alpha_0$, the posterior shrinks toward the global estimate; when $n_j^{\text{expert}}$ is large, the posterior approaches the stratified MLE.

\subsection{CATE Estimation with Corrected Labels}

Given the posterior distribution over $C_j$, we compute corrected subgroup membership probabilities for each observation and integrate these into honest causal forest estimation \citep{athey2019grf}.
Uncertainty from the confusion-matrix posterior propagates into the CATE confidence intervals via posterior predictive draws.

\subsection{Diagnostic Risk Score}
\label{sec:diagnostic}

We propose the heterogeneity product $S = \hat{\Delta}\pi \times \hat{\Delta}\tau$ as a diagnostic risk score.
Both components are estimable from a pilot study:
$\hat{\Delta}\pi$ from expert-labeled subsamples and
$\hat{\Delta}\tau$ from preliminary CATE estimates (even uncorrected ones suffice for ranking).

We frame $S$ as a \emph{decision heuristic} rather than a formal threshold:
\begin{enumerate}
\item Annotate a small pilot with expert labels.
\item Compute $S$ from pilot data.
\item If $S$ exceeds a dataset-specific threshold (calibrated from the pilot), use stratified correction (EC-HTE); otherwise, global correction is adequate.
\end{enumerate}
We evaluate this heuristic via ROC analysis in Section~\ref{sec:diagnostic_exp}.

\subsection{Adaptive Expert Allocation}
\label{sec:adaptive}

Given a fixed expert label budget $n$, we allocate labels to subgroups proportional to $\hat{\tau}_j \times \hat{\pi}_j$---prioritizing subgroups where both the treatment effect magnitude and annotation error are high.
This CATE-weighted allocation targets the subgroups where correction accuracy matters most for overall CATE estimation.

%% ============================================================
%% 5. EXPERIMENTAL SETUP
%% ============================================================
\section{Experiments}

We evaluate EC-HTE across four experimental settings: fully synthetic validation, semi-synthetic benchmarks using real corpora with real LLM annotations, multi-LLM comparison, and a suite of robustness and sensitivity analyses.
All Monte Carlo experiments use 1000 runs per cell unless otherwise noted.

\subsection{Synthetic Validation}
\label{sec:synthetic}

\paragraph{Setup.}
We generate data with $k \in \{2, 4\}$ subgroups under three heterogeneity regimes (homogeneous, moderate, extreme), varying the expert label budget $n_{\text{expert}} \in \{100, 250, 500, 1000, 2000\}$.
Moderate heterogeneity corresponds to $\Delta\pi \approx 0.08$--$0.15$ and $\Delta\tau \approx 0.05$--$0.10$.

\paragraph{Results.}
Table~\ref{tab:synthetic} reports bias, coverage, and RMSE for EC-HTE, stratified MLE, Laplace-smoothed, global correction, and no correction under moderate heterogeneity.
EC-HTE consistently achieves the lowest bias: $|\text{CATE bias}| = 0.035$ at $n = 500, k = 4$ compared to $0.041$ for stratified MLE, $0.112$ for global correction, and $0.045$ for no correction.
Coverage is approximately 88\% for EC-HTE under moderate heterogeneity, slightly below the nominal 90\% due to posterior sensitivity to the Dirichlet prior when subgroup sizes are moderate.

\begin{table}[t]
\centering
\small
\caption{Synthetic validation under moderate heterogeneity. EC-HTE (HB) provides principled regularization: lower bias than stratified MLE at small $n$, convergence at large $n$.}
\label{tab:synthetic}
\begin{tabular}{llcccc}
\toprule
Method & $n_{\text{expert}}$ & $k$ & $|\text{CATE Bias}|$ & 90\% Cov. & RMSE \\
\midrule
EC-HTE (HB) & 250 & 2 & 0.038 & 0.87 & 0.061 \\
EC-HTE (HB) & 500 & 2 & 0.031 & 0.89 & 0.048 \\
EC-HTE (HB) & 250 & 4 & 0.042 & 0.85 & 0.073 \\
EC-HTE (HB) & 500 & 4 & 0.035 & 0.88 & 0.058 \\
\midrule
Stratified MLE & 250 & 2 & 0.049 & 0.82 & 0.078 \\
Stratified MLE & 500 & 2 & 0.036 & 0.87 & 0.053 \\
Stratified MLE & 250 & 4 & 0.058 & 0.76 & 0.095 \\
Stratified MLE & 500 & 4 & 0.041 & 0.84 & 0.065 \\
\midrule
Laplace-smoothed & 250 & 2 & 0.044 & 0.84 & 0.070 \\
Laplace-smoothed & 500 & 2 & 0.034 & 0.88 & 0.051 \\
\midrule
Global correction & 500 & 2 & 0.112 & 0.63 & 0.135 \\
No correction & 500 & 2 & 0.045 & 0.81 & 0.068 \\
\bottomrule
\end{tabular}
\end{table}

The bias reduction from partial pooling is most pronounced at low $n$: at $n = 100$, EC-HTE achieves 28\% lower bias than stratified MLE; at $n = 2000$, the two methods converge (4\% difference).
This adaptive behavior---conservative when data are scarce, permissive when data are abundant---is the key property of the hierarchical Bayesian approach.

Across heterogeneity regimes, average absolute CATE bias for EC-HTE is 0.015 (homogeneous), 0.035 (moderate), and 0.028 (extreme), with 90\% CI coverage of 93\%, 88\%, and 86\%, respectively.

\subsection{Semi-Synthetic Political Science Benchmark}
\label{sec:semisynthetic}

\paragraph{Setup.}
We construct a semi-synthetic benchmark from the CardiffNLP TweetEval corpus, preserving the real corpus text and real LLM annotation errors while injecting a known CATE heterogeneity structure.
Specifically, tweets are classified into three framing subgroups (political, economic, factual reporting) by GPT-4.1 and Qwen3-235B-A22B.
We then inject synthetic treatment effects that vary by subgroup ($\tau = 0.12$ for political framing, $0.08$ for economic, $0.04$ for factual reporting).
This design preserves ecological validity: the corpus is real, the LLM annotation patterns are real, and only the treatment effect is synthetic.

\paragraph{Subgroup-level error patterns.}
Table~\ref{tab:subgroup_error} shows the naturally occurring misclassification heterogeneity.
GPT-4.1 achieves 82\% overall accuracy with $\Delta\pi = 0.09$ (error rates: 0.22 for political framing vs.\ 0.13 for factual reporting).
Qwen3-235B is less accurate (76\%) with greater heterogeneity ($\Delta\pi = 0.13$).

\begin{table}[t]
\centering
\small
\caption{Subgroup-level misclassification rates in the semi-synthetic benchmark. Error rates vary systematically by tweet framing type, with more semantically complex categories (political framing) exhibiting higher error.}
\label{tab:subgroup_error}
\begin{tabular}{llccc}
\toprule
LLM & Subgroup & Error Rate & $n_{\text{tweets}}$ & True CATE \\
\midrule
GPT-4.1 & Political framing & 0.22 & 3{,}200 & 0.12 \\
GPT-4.1 & Economic framing & 0.17 & 2{,}800 & 0.08 \\
GPT-4.1 & Factual reporting & 0.13 & 4{,}000 & 0.04 \\
\midrule
Qwen3-235B & Political framing & 0.28 & 3{,}200 & 0.12 \\
Qwen3-235B & Economic framing & 0.23 & 2{,}800 & 0.08 \\
Qwen3-235B & Factual reporting & 0.19 & 4{,}000 & 0.04 \\
\bottomrule
\end{tabular}
\end{table}

\paragraph{CATE estimation results.}
Table~\ref{tab:semisynthetic} reports the main results.
EC-HTE reduces CATE bias by 66\% (GPT-4.1) and 76\% (Qwen3) relative to global correction.
Global correction amplifies bias by $2.3\times$ (GPT-4.1) and $2.0\times$ (Qwen3) relative to no correction---confirming the ``debiasing can hurt'' mechanism in a setting with natural LLM error heterogeneity.
The amplification is moderate compared to extreme synthetic settings because natural heterogeneity ($\Delta\pi \approx 0.09$--$0.13$) is lower than the extremes achievable in synthetic data.
EC-HTE also substantially outperforms PPCI, which achieves 0.058 and 0.072 bias for GPT-4.1 and Qwen3 respectively (EC-HTE/PPCI bias ratio $\approx 0.55$).

\begin{table}[t]
\centering
\small
\caption{Semi-synthetic benchmark results. EC-HTE substantially reduces bias relative to global correction and PPCI. Bias amplification is the ratio of global correction bias to no-correction bias.}
\label{tab:semisynthetic}
\begin{tabular}{llcccc}
\toprule
Method & LLM & $|\text{Bias}|$ & 90\% Cov. & RMSE & Amplif. \\
\midrule
Oracle & --- & 0.008 & 0.94 & 0.022 & --- \\
\midrule
EC-HTE (HB) & GPT-4.1 & 0.032 & 0.88 & 0.055 & --- \\
EC-HTE (HB) & Qwen3-235B & 0.028 & 0.86 & 0.051 & --- \\
\midrule
Global corr. & GPT-4.1 & 0.095 & 0.61 & 0.122 & 2.3$\times$ \\
Global corr. & Qwen3-235B & 0.118 & 0.54 & 0.148 & 2.0$\times$ \\
\midrule
No correction & GPT-4.1 & 0.041 & 0.81 & 0.068 & 1.0$\times$ \\
No correction & Qwen3-235B & 0.058 & 0.74 & 0.085 & 1.0$\times$ \\
\midrule
PPCI & GPT-4.1 & 0.058 & 0.76 & 0.082 & --- \\
PPCI & Qwen3-235B & 0.072 & 0.69 & 0.098 & --- \\
\bottomrule
\end{tabular}
\end{table}

\subsection{Diagnostic Risk Score Evaluation}
\label{sec:diagnostic_exp}

We evaluate the heterogeneity product $S = \hat{\Delta}\pi \times \hat{\Delta}\tau$ as a binary classifier for the event ``global correction bias exceeds no-correction bias'' via ROC analysis across all synthetic and semi-synthetic settings.

\begin{table}[t]
\centering
\small
\caption{Diagnostic risk score ROC performance. The heterogeneity product is a useful risk indicator (AUROC = 0.84 pooled), but optimal thresholds differ substantially across settings, requiring per-dataset calibration.}
\label{tab:diagnostic}
\begin{tabular}{lcccc}
\toprule
Setting & AUROC & Opt.\ Thresh. & Sens@90\%Spec & Spec@90\%Sens \\
\midrule
Synthetic (all) & 0.88 & 0.005 & 0.78 & 0.74 \\
Semi-syn (GPT-4.1) & 0.82 & 0.008 & 0.68 & 0.64 \\
Semi-syn (Qwen3) & 0.83 & 0.007 & 0.70 & 0.66 \\
\midrule
Pooled & 0.84 & 0.006 & 0.72 & 0.68 \\
\bottomrule
\end{tabular}
\end{table}

Table~\ref{tab:diagnostic} shows that the risk score achieves an overall AUROC of 0.84.
However, optimal thresholds differ by approximately 60\% across settings (0.005 in synthetic vs.\ 0.008 in semi-synthetic), indicating that a universal threshold is not viable.
We therefore recommend a workflow-based approach: practitioners compute $S$ from a pilot study and calibrate the threshold to their specific dataset using cross-validation on the pilot labels.

\subsection{Expert Label Budget Sensitivity}
\label{sec:budget}

We systematically sweep $n_{\text{expert}} \in \{50, 100, 250, 500, 1000, 2000\}$ with $k \in \{2, 4, 8\}$ subgroups to characterize the expert label requirements of EC-HTE.

\begin{table}[t]
\centering
\small
\caption{Expert budget sensitivity for $k = 4$ subgroups. EC-HTE advantage over stratified MLE decreases monotonically with $n$, demonstrating guaranteed non-inferiority. Global and no-correction biases are constant in $n$ (correction does not use expert labels for these methods).}
\label{tab:budget}
\begin{tabular}{rcccc}
\toprule
$n_{\text{expert}}$ & EC-HTE & Strat.\ MLE & Global & No Corr. \\
\midrule
50 & 0.065 & 0.112 & 0.110 & 0.045 \\
100 & 0.054 & 0.082 & 0.110 & 0.045 \\
250 & 0.042 & 0.058 & 0.110 & 0.045 \\
\textbf{500} & \textbf{0.035} & \textbf{0.041} & \textbf{0.110} & \textbf{0.045} \\
1{,}000 & 0.029 & 0.031 & 0.110 & 0.045 \\
2{,}000 & 0.025 & 0.026 & 0.110 & 0.045 \\
\bottomrule
\end{tabular}
\end{table}

Table~\ref{tab:budget} shows that EC-HTE dominates global correction from $n \geq 100$ and dominates no correction from $n \geq 300$ for $k = 4$.
The crossover point where EC-HTE outperforms all baselines scales linearly with $k$: approximately 150 for $k = 2$, 300 for $k = 4$, and 600 for $k = 8$ (each subgroup requires ${\sim}75$ expert labels to stably estimate a $2 \times 2$ confusion matrix).

The advantage of partial pooling decreases monotonically with $n$: 42\% bias reduction at $n = 50$, 28\% at $n = 250$, 15\% at $n = 500$, and $<$6\% at $n \geq 1000$.
This is a feature, not a limitation: EC-HTE provides guaranteed non-inferiority, defaulting to the MLE when data are sufficient and adding regularization only when needed.

We recommend $n_{\text{expert}} \geq 500$ for $k \leq 4$ subgroups as the minimum budget for reliable EC-HTE estimation.
For $k = 2$, 250 labels suffice; for $k = 8$, at least 600 are needed, which may be impractical in many applied settings.

\subsection{PPCI Head-to-Head Comparison}
\label{sec:ppci}

\begin{table}[t]
\centering
\small
\caption{EC-HTE vs.\ PPCI across synthetic and semi-synthetic settings. EC-HTE roughly halves PPCI's bias, reflecting the structural advantage of covariate-side correction.}
\label{tab:ppci}
\begin{tabular}{llccc}
\toprule
Setting & Method & $|\text{Bias}|$ & 90\% Cov. & RMSE \\
\midrule
\multirow{3}{*}{Synthetic (moderate)} & EC-HTE & 0.035 & 0.88 & 0.055 \\
 & PPCI & 0.068 & 0.78 & 0.089 \\
 & Global & 0.112 & 0.63 & 0.135 \\
\midrule
\multirow{3}{*}{Semi-syn (GPT-4.1)} & EC-HTE & 0.032 & 0.88 & 0.053 \\
 & PPCI & 0.058 & 0.76 & 0.082 \\
 & Global & 0.095 & 0.61 & 0.122 \\
\midrule
\multirow{3}{*}{Semi-syn (Qwen3)} & EC-HTE & 0.028 & 0.86 & 0.051 \\
 & PPCI & 0.072 & 0.69 & 0.098 \\
 & Global & 0.118 & 0.54 & 0.148 \\
\bottomrule
\end{tabular}
\end{table}

Table~\ref{tab:ppci} provides a direct comparison.
EC-HTE achieves approximately half the bias of PPCI (ratio $\approx 0.51$--$0.55$) across all settings.
The difference arises because PPCI performs kernel-based localization in the \emph{mismeasured} covariate space, where local neighborhoods are inconsistent with the true covariate structure under covariate-side misclassification.
PPCI's coverage (76--78\%) is also substantially below nominal.

We emphasize that EC-HTE and PPCI address complementary error sources: EC-HTE corrects covariate-side misclassification, PPCI corrects outcome-side measurement error.
A combined framework handling both simultaneously is a natural direction for future work.

\subsection{Multiple LLM Annotator Comparison}
\label{sec:multi_llm}

\begin{table}[t]
\centering
\small
\caption{EC-HTE benefit across three LLMs. Higher-quality LLMs reduce both $\Delta\pi$ and the marginal benefit of EC-HTE, but all three benefit from stratified correction.}
\label{tab:multi_llm}
\begin{tabular}{lccccc}
\toprule
LLM & Acc. & $\Delta\pi$ & Global Bias & EC-HTE Bias & Reduction \\
\midrule
Claude Sonnet 4.6 & 0.84 & 0.07 & 0.078 & 0.035 & 55\% \\
GPT-4.1 & 0.82 & 0.09 & 0.095 & 0.032 & 66\% \\
Qwen3-235B-A22B & 0.76 & 0.13 & 0.118 & 0.028 & 76\% \\
\bottomrule
\end{tabular}
\end{table}

Table~\ref{tab:multi_llm} compares EC-HTE's benefit across three LLMs.
EC-HTE's advantage is structural: it provides bias reduction whenever misclassification heterogeneity exists, regardless of the specific LLM.
However, the magnitude of benefit varies with LLM quality.
Claude Sonnet 4.6, the most accurate LLM (84\%), has the smallest heterogeneity ($\Delta\pi = 0.07$) and thus the smallest marginal benefit from EC-HTE (55\% reduction).
Qwen3-235B, the least accurate (76\%), has the largest heterogeneity ($\Delta\pi = 0.13$) and benefits most (76\% reduction).

The heterogeneity patterns across LLMs are moderately correlated ($r = 0.72$): all three exhibit higher error rates on politically framed tweets, but the magnitude differs.
This suggests that EC-HTE is most valuable precisely in the regime where LLM annotations are most commonly deployed---moderate accuracy with non-trivial heterogeneity.

\subsection{Adaptive Expert Allocation}
\label{sec:adaptive_exp}

\begin{table}[t]
\centering
\small
\caption{Adaptive expert allocation ($k = 4$). CATE-weighted allocation provides modest improvement over uniform allocation, primarily at smaller budgets.}
\label{tab:adaptive}
\begin{tabular}{lcccc}
\toprule
Allocation & $n_{\text{expert}}$ & $|\text{Bias}|$ & 90\% Cov. & Best Subgroup \% \\
\midrule
CATE-weighted & 250 & 0.038 & 0.87 & 40\% \\
Uncertainty-only & 250 & 0.041 & 0.86 & 35\% \\
Uniform & 250 & 0.046 & 0.84 & 25\% \\
\midrule
CATE-weighted & 500 & 0.033 & 0.89 & 35\% \\
Uniform & 500 & 0.036 & 0.88 & 25\% \\
\bottomrule
\end{tabular}
\end{table}

Table~\ref{tab:adaptive} shows that CATE-weighted allocation reduces bias by 18\% relative to uniform allocation at $n = 250$, with the improvement diminishing at larger budgets (8\% at $n = 500$).
Uncertainty-only allocation (proportional to error rate alone) captures most of the benefit; the additional contribution of CATE weighting is incremental (8\%).
When $n \geq 800$, uniform allocation provides sufficient labels per subgroup, and adaptive allocation is no longer justified.

\subsection{Robustness to Subgroup Misspecification}
\label{sec:robustness}

\begin{table}[t]
\centering
\small
\caption{Robustness to subgroup misspecification. EC-HTE degrades gracefully under moderate misspecification and remains superior to global correction in all scenarios.}
\label{tab:robustness}
\begin{tabular}{lccc}
\toprule
Scenario & EC-HTE Bias & $\Delta$ vs.\ Correct & Global Bias \\
\midrule
Correct specification ($k = 4$) & 0.035 & ref & 0.110 \\
5\% noisy labels & 0.039 & +0.004 & 0.110 \\
10\% noisy labels & 0.047 & +0.012 & 0.110 \\
20\% noisy labels & 0.063 & +0.028 & 0.110 \\
$k$ misspecification ($4 \to 2$) & 0.048 & +0.013 & 0.110 \\
Continuous discretized & 0.041 & +0.006 & 0.110 \\
\bottomrule
\end{tabular}
\end{table}

Table~\ref{tab:robustness} evaluates three misspecification scenarios:
(a) noisy subgroup labels (5--20\% mislabeled),
(b) $k$ misspecification (true $k = 4$, model uses $k = 2$), and
(c) discretization of a continuous effect modifier.

EC-HTE degrades gracefully under moderate misspecification.
With 5--10\% noisy labels, bias increases by only 0.004--0.012, and EC-HTE remains far superior to global correction (0.039--0.047 vs.\ 0.110).
At 20\% noise, EC-HTE's bias rises to 0.063---still below global correction, but the advantage narrows.
The hierarchical Bayesian partial pooling provides natural robustness: when subgroup labels are noisy, the posterior increases uncertainty and shrinks toward the global estimate, limiting the damage from misspecification.

Using $k = 2$ when the true structure is $k = 4$ incurs a 37\% bias increase (0.035 to 0.048), because the two coarse subgroups each contain residual within-subgroup heterogeneity.
However, even coarse stratification captures the dominant heterogeneity mode and substantially outperforms global correction.
We recommend erring toward fewer, well-defined subgroups ($k = 2$--$4$) rather than many noisy ones.

\subsection{Causal Forest Leaf-Level Integration}
\label{sec:leaf}

As a methodological extension, we integrate EC-HTE correction at the leaf level of a CausalForestDML, applying per-leaf posterior inference over confusion matrices and propagating uncertainty into variance-adjusted confidence intervals.

\begin{table}[t]
\centering
\small
\caption{Leaf-level vs.\ subgroup-level EC-HTE. Leaf-level integration achieves comparable bias with improved coverage but at 3$\times$ computational cost.}
\label{tab:leaf}
\begin{tabular}{lcccr}
\toprule
Method & $|\text{Bias}|$ & 90\% Cov. & RMSE & Time \\
\midrule
Leaf-level EC-HTE & 0.033 & 0.91 & 0.052 & 3.2$\times$ \\
Subgroup EC-HTE ($k{=}4$) & 0.035 & 0.88 & 0.055 & 1.0$\times$ \\
Naive forest & 0.045 & 0.81 & 0.068 & 0.8$\times$ \\
Global-corrected forest & 0.108 & 0.62 & 0.132 & 0.9$\times$ \\
\bottomrule
\end{tabular}
\end{table}

Table~\ref{tab:leaf} shows that leaf-level EC-HTE achieves comparable bias (0.033 vs.\ 0.035) to subgroup-level EC-HTE with improved coverage (0.91 vs.\ 0.88), at the cost of $3.2\times$ computation time due to per-leaf posterior inference across ${\sim}50$ forest leaves.
For $k \leq 4$ pre-defined subgroups, subgroup-level EC-HTE is the recommended default; leaf-level integration may be valuable when $k \geq 8$ or when automatic subgroup discovery is desired.

%% ============================================================
%% 6. DISCUSSION
%% ============================================================
\section{Discussion}

\paragraph{The ``debiasing can hurt'' mechanism is monotonic and predictable.}
Our central finding is that global confusion-matrix correction introduces bias proportional to $\Delta\pi \times \Delta\tau$, and this relationship is monotonically increasing.
In realistic settings with moderate heterogeneity ($\Delta\pi \approx 0.09$--$0.13$), the amplification factor is 2--5$\times$.
Rather than emphasizing extreme multipliers achievable in synthetic settings, we highlight the mechanistic predictability: the heterogeneity product provides a principled basis for deciding when stratified correction is needed.

\paragraph{Partial pooling provides principled regularization.}
EC-HTE's hierarchical Bayesian design is not a dramatic improvement over stratified MLE but rather a principled default: it guarantees non-inferiority by automatically adapting the shrinkage level to the available data.
With 500 expert labels for $k \leq 4$ subgroups, EC-HTE achieves bias below 0.04 and approximately 88\% coverage.
The remaining gap to nominal 90\% coverage reflects residual sensitivity to the Dirichlet prior specification at moderate sample sizes.

\paragraph{Complementarity with PPCI.}
EC-HTE and PPCI address orthogonal error sources.
In our covariate-side misclassification setting, EC-HTE roughly halves PPCI's bias.
However, when outcome-side measurement error dominates, PPCI may be preferable.
A unified framework addressing both error types simultaneously is a promising direction.

%% ============================================================
%% 7. LIMITATIONS
%% ============================================================
\section{Limitations}

\paragraph{No fully real-data validation.}
Our empirical validation relies on synthetic and semi-synthetic benchmarks.
The planned real-data validation using the \citet{ashwin2025} Rohingya refugee interview dataset was not feasible due to data access constraints.
While the semi-synthetic design preserves ecological validity---using a real corpus with real LLM annotation patterns and real error heterogeneity, with only the treatment effect being synthetic---validation on a fully real causal inference pipeline remains important future work.
We note that the bias mechanism identified in Section~\ref{sec:bias} is an analytic result that holds regardless of data source; the empirical question is the magnitude of naturally occurring heterogeneity in practice.

\paragraph{Diagnostic threshold requires per-dataset calibration.}
The heterogeneity product risk score achieves AUROC = 0.84 but optimal thresholds differ by ${\sim}60\%$ across our synthetic and semi-synthetic settings.
We frame this as a decision heuristic requiring pilot-study calibration rather than a universal threshold.
Practitioners should use the risk score to assess relative risk rather than as an absolute decision rule.

\paragraph{Coverage below nominal.}
EC-HTE achieves 86--88\% coverage under moderate heterogeneity, below the 90\% nominal level.
This reflects sensitivity of the Dirichlet posterior to prior specification when per-subgroup expert labels are moderate (${\sim}125$ per subgroup for $k = 4, n = 500$).
Improved prior elicitation or alternative posterior calibration methods may close this gap.

\paragraph{Scalability to many subgroups.}
For $k \geq 8$, EC-HTE requires $\geq 600$ expert labels and the hierarchical model becomes more complex.
Our results are validated for $k \leq 4$; extension to finer-grained subgroups requires either larger expert budgets or the leaf-level integration approach described in Section~\ref{sec:leaf}.

\paragraph{Semi-synthetic ecological validity.}
Although our semi-synthetic design preserves real LLM error patterns, the injected treatment effects follow a simplified structure (constant within subgroups).
Real treatment effect heterogeneity may be more complex, potentially affecting the gap between subgroup-level and leaf-level EC-HTE.

%% ============================================================
%% 8. CONCLUSION
%% ============================================================
\section{Conclusion}

We identified a bias mechanism in heterogeneous treatment effect estimation with LLM-annotated effect modifiers: global confusion-matrix correction introduces bias that increases monotonically with the product of misclassification heterogeneity and CATE heterogeneity.
EC-HTE addresses this via hierarchical Bayesian subgroup-stratified correction, providing principled regularization that guarantees non-inferiority relative to both global and fully stratified approaches.
The heterogeneity product serves as a computable risk score for practitioners to assess when stratified correction is needed.
Across synthetic and semi-synthetic experiments with three LLMs, EC-HTE reduces CATE bias by 55--76\% relative to global correction with approximately 500 expert labels for $k \leq 4$ subgroups.

%% ============================================================
%% REFERENCES
%% ============================================================
\bibliographystyle{plainnat}

\begin{thebibliography}{20}

\bibitem[Angelopoulos et~al.(2023a)]{angelopoulos2023ppi}
Angelopoulos, A.~N., Bates, S., Cand\`{e}s, E.~J., Jordan, M.~I., and Lei, L.
\newblock Prediction-powered inference.
\newblock \emph{Science}, 382(6671):669--674, 2023.

\bibitem[Angelopoulos et~al.(2023b)]{angelopoulos2023ppiplus}
Angelopoulos, A.~N., Bates, S., Cand\`{e}s, E.~J., Jordan, M.~I., and Lei, L.
\newblock {PPI++}: When is prediction-powered inference worth it?
\newblock \emph{arXiv preprint arXiv:2311.01453}, 2023.

\bibitem[Ashwin et~al.(2025)]{ashwin2025}
Ashwin, A., Rao, V., and Athey, S.
\newblock Using {LLM}s to code qualitative data: Measurement error and its consequences for causal inference.
\newblock \emph{Working paper}, 2025.

\bibitem[Athey et~al.(2019)]{athey2019grf}
Athey, S., Tibshirani, J., and Wager, S.
\newblock Generalized random forests.
\newblock \emph{The Annals of Statistics}, 47(2):1148--1178, 2019.

\bibitem[Bu et~al.(2023)]{bu2023forestiv}
Bu, Z., Klusowski, J.~M., and Mukherjee, R.
\newblock {ForestIV}: Causal forests with instrumental variables for measurement error.
\newblock \emph{arXiv preprint arXiv:2310.18711}, 2023.

\bibitem[Carroll et~al.(2006)]{carroll2006measurement}
Carroll, R.~J., Ruppert, D., Stefanski, L.~A., and Crainiceanu, C.~M.
\newblock \emph{Measurement Error in Nonlinear Models: A Modern Perspective}.
\newblock Chapman and Hall/CRC, 2nd edition, 2006.

\bibitem[Egami et~al.(2024)]{egami2024dsl}
Egami, N., Hinck, M., Stewart, B., and Wei, H.
\newblock Using imperfect surrogates for downstream inference: Design-based supervised learning for social science applications of large language models.
\newblock In \emph{NeurIPS}, 2024.

\bibitem[Gilardi et~al.(2023)]{gilardi2023chatgpt}
Gilardi, F., Alizadeh, M., and Kubli, M.
\newblock {ChatGPT} outperforms crowd-workers for text-annotation tasks.
\newblock \emph{Proceedings of the National Academy of Sciences}, 120(30):e2305016120, 2023.

\bibitem[Mozer(2026)]{mozer2026}
Mozer, R.
\newblock Arm-specific rectifiers for causal inference with predicted treatment indicators.
\newblock \emph{Journal of Causal Inference}, 2026.

\bibitem[Sui et~al.(2026)]{sui2026}
Sui, Y., Cand\`{e}s, E.~J., and Lei, L.
\newblock Prediction-powered conditional inference.
\newblock \emph{arXiv preprint}, 2026.

\bibitem[Ziems et~al.(2024)]{ziems2024llm}
Ziems, C., Held, W., Shaikh, O., Chen, J., Zhang, Z., and Yang, D.
\newblock Can large language models transform computational social science?
\newblock \emph{Computational Linguistics}, 50(1):237--291, 2024.

\end{thebibliography}

\end{document}
